\begin{proposition}[A projective example with twelve length-one contributions]\hypertarget{tag-01BL}{}\label{cdp:st:B2}
\StatementRecord{cdp:st:B2}{proved}{01BL}{\SeriesRef{cdp:eq:B.6}, \SeriesRef{cdp:eq:B.5}, \SeriesRef{cdp:eq:B.7}, \SeriesRef{cdp:st:B1}}

Let $h:E\to\PP^1_t$ be the Weierstrass surface defined on the finite chart by
\[
 y^2z=x^3+(t^4+1)xz^2+t^6z^3,                          \tag{B.3}\hypertarget{tag-008I}{}\label{cdp:eq:B.3}
\]
and at $s=1/t$ by $x_\infty=s^2x$, $y_\infty=s^3y$, $z_\infty=z$. Let $O=[0:1:0]$ and $P=(0,t^3,1)$ be its two sections and put
\[
 M_E=\OO_E(P-O),\quad Y=E\times\PP^1,\quad
 g=h\circ\operatorname{pr}_E,\quad
 L=K_Y\otimes\operatorname{pr}_E^*M_E^{-1},\quad
 \mathcal E=T_Y\otimes L.
\]
The surface $E$ is smooth and projective. Its singular fibres are twelve irreducible nodal cubics, at the distinct roots $t_i$ of
\[
 D(t)=31t^{12}+12t^8+12t^4+4.
\]
For every such fibre $S_i=g^{-1}(t_i)$, with $A=(K_Y\otimes L^{-1})|_{S_i}$, we have
\[
H^0(S_i,N^*_{S_i/Y}\otimes A)=0,
\qquad H^0(S_i,\widetilde\Omega^1_{S_i}\otimes A)=0,
\]
\[
 (h^0,h^1,h^2)(S_i,\mathcal E|_{S_i})=(0,1,1).
\]
As coherent sheaves on the base, with noncanonical identifications of their one-dimensional fibres,
\[
 R^qg_*\mathcal E=0\ (q\ne2),\qquad
 R^2g_*\mathcal E=\bigoplus_{i=1}^{12}\C_{t_i}.        \tag{B.4}\hypertarget{tag-008J}{}\label{cdp:eq:B.4}
\]
In particular $(h^0,h^1,h^2,h^3)(Y,\mathcal E)=(0,0,12,0)$.
\end{proposition}



\begin{proof}

\begin{enumerate}[label=\arabic*.,leftmargin=*]
\item The transition rule realizes \eqref{cdp:eq:B.3} as a cubic in a projective plane bundle over $\PP^1$, with Weierstrass coefficients in $\OO(4)$ and $\OO(6)$; hence its total space is projective. Its discriminant is $-16D(t)$. Exact polynomial division gives
\[
 \gcd(D,D')=1,\qquad \gcd(D,t^4+1)=1.
\]
The leading coefficient $31$ is its nonzero value in the chart at infinity. Each discriminant root therefore has nonzero $c_4=-48(t^4+1)$ and is a nodal cubic. A simple discriminant zero is transverse to the nodal discriminant and makes the surface smooth at that node. Alternatively the Gr\"obner basis of the affine total-space singularity ideal $(F,F_x,F_y,F_t)$ is $(1)$. At the zero section smoothness follows from $F_z=1$ in the chart $y=1$. At infinity the coefficients are $a=1+s^4,b=1$, with discriminant $-16\cdot31$ at $s=0$.

\item The sections $P,O$ are disjoint and meet the smooth locus of every fibre. For $P$, a singular point would require both $a(t)=0$ and $t=0$, which is impossible. Each fibre is an integral arithmetic-genus-one curve with trivial dualizing sheaf. The line bundle $M_E$ has degree zero on every fibre and is nontrivial there. On a smooth fibre, a trivialization would give a meromorphic function with one simple pole and one simple zero, hence a degree-one map to $\PP^1$, contradicting genus one. On a nodal fibre the normalization calculation in \eqref{cdp:eq:B.2} says that two distinct smooth points determine a nontrivial degree-zero line bundle. Thus $H^q(E_t,M_E^{\pm1}|_{E_t})=0$ for both $q=0,1$: a nonzero section of a degree-zero line bundle on an integral curve is a trivialization, and Serre duality gives the assertion in degree one. Grauert's theorem, applied to the proper flat map $h$, yields
\[
 Rh_*M_E=Rh_*M_E^{-1}=0.                               \tag{B.5}\hypertarget{tag-008K}{}\label{cdp:eq:B.5}
\]

\item Let $Z\subset E$ be the reduced set of the twelve critical points. The relative dualizing sheaf is $\omega_{E/\PP^1}=h^*\OO(1)$. Indeed the residue frame $dx/(2y)$ has this transition under $x_\infty=s^2x,y_\infty=s^3y$, and extends as a dualizing frame at a node. It follows that $K_E=h^*\OO(-1)$. There is an exact sequence of coherent sheaves on $E$
\[
0\longrightarrow h^*\Omega^1_{\PP^1}
 \xrightarrow{dh}\Omega^1_E
 \longrightarrow\mathcal I_Z\otimes\omega_{E/\PP^1}
 \longrightarrow0.                                   \tag{B.6}\hypertarget{tag-008L}{}\label{cdp:eq:B.6}
\]
Away from $Z$ this is the ordinary relative cotangent sequence. At a node choose coordinates $vw=t$. The map from $\Omega^1_E$ to the relative dualizing line sends $dv\mapsto v\omega$, $dw\mapsto-w\omega$; its image is $(v,w)\omega$ and its kernel is generated by $w\,dv+v\,dw=dt$. This proves \eqref{cdp:eq:B.6} on every stalk, with the critical scheme reduced.

\item Tensor \eqref{cdp:eq:B.6} by $M_E^{-1}$. By \eqref{cdp:eq:B.5}, projection formula, and the sequence for the ideal of $Z$, its derived direct image has only degree one, namely
\[
 R^1h_*(\Omega^1_E\otimes M_E^{-1})
 =h_*(\omega_{E/\PP^1}\otimes M_E^{-1}|_Z),
 \qquad h_*(\Omega^1_E\otimes M_E^{-1})=0.             \tag{B.7}\hypertarget{tag-008M}{}\label{cdp:eq:B.7}
\]
The right side is the sum of twelve reduced skyscrapers. Each summand in \eqref{cdp:eq:B.7} has length one.

\item Contraction on the smooth surface gives the canonical isomorphism $T_E\otimes K_E\simeq\Omega^1_E$. Consequently the product tangent decomposition identifies
\[
\mathcal E=(\Omega^1_E\otimes M_E^{-1})\boxtimes\OO(-2)
 \ \oplus\ (K_E\otimes M_E^{-1})\boxtimes\OO.
\]
The second summand has zero derived direct image by \eqref{cdp:eq:B.5}. The first has only the tensor product of degree-one terms: \eqref{cdp:eq:B.7} and $H^1(\PP^1,\OO(-2))=\C$. The relative coherent K\"unneth comparison gives \eqref{cdp:eq:B.4}. The Leray spectral sequence over a curve then gives the four stated global dimensions.

\item On a singular fibre $S_i=E_{t_i}\times\PP^1$, the restriction of $M_E$ is a nontrivial $A_\lambda$. Statement \ref{cdp:st:B1} gives the two vanishings and ambient dimensions $(1,1,0)$. Adjunction for the Cartier fibre gives $\omega_{S_i}=K_Y|_{S_i}$, because its normal line is trivial. Serre duality for this compact Gorenstein surface identifies, for $q=0,1,2$,
\[
 H^q(S_i,\mathcal E|_{S_i})^\vee
 =H^{2-q}(S_i,\Omega_Y^1|_{S_i}\otimes A).
\]
The fibre value $h^1=1$ with $R^1g_*\mathcal E=0$ is the degree-one Tor term obtained by tensoring a reduced skyscraper in degree two with the residue field.

\end{enumerate}
\end{proof}